\documentclass[12pt]{article}


\usepackage{fullpage}
\usepackage{amsmath}
\usepackage{amssymb}
\usepackage{color}
\date{}

\newcommand{\mys}[2]{\displaystyle\sum_{#1}^{#2}}

\usepackage{tikz}

\newif\ifans
\anstrue


\begin{document}

\begin{center}
  \textbf{COMS 3110: Homework 2}\\
  \textbf{Due: October $11^{th}$, 11:59pm}\\
  \textbf{Total Points: 50}
\end{center}


\paragraph{Late submission policy.\ }
An assignment that is submitted one day late will get a penalty of $10\%$. An assignment that is submitted two days late will get a penalty of $20\%$. Assignments submitted after two days will not be graded and will get no points. For example, if an assignment is due on Friday by midnight, a submission on Saturday will be penalized by $10\%$, and a submission on Sunday will be penalized by $20\%$. A submission on Monday will get no points. 

\paragraph{Submission format.\ }
Homework solutions will have to be typed. You can use word, LaTeX, or any other type-setting tool to type your solution. Your submission file should be in pdf format. Do \textbf{NOT} submit a photocopy of handwritten homework except for diagrams that can be hand-drawn and scanned. We reserve the right \textbf{NOT} to grade homework that does not follow the formatting requirements.
Name your submission
file: \texttt{<Your-net-id>-3110-hw2.pdf}. For instance, if your netid
is \texttt{asterix}, then your submission file will be named
\texttt{asterix-3110-hw2.pdf}.
Each student must hand in their own assignment. If you discussed the homework or solutions with others, a list of collaborators must be included with each submission. Each of the collaborators has to write the solutions in their own words (copies are not allowed).

\paragraph{General Requirements}
\begin{itemize}
    \item When proofs are required, do your best to make them both clear and rigorous. Even when proofs are not required, you should justify your answers and explain your work.
    \item When asked to present a construction, you should show the correctness of the construction.
\end{itemize}

\paragraph{Some Useful (in)equalities}
\begin{itemize}
    \item $\sum_{i=1}^n i= \frac{n(n+1)}{2}$
    \item $\sum_{i=1}^ni^2 = \frac{n(n+1)(2n+1)}{6}$
    \item $2^{\log_2n}=n$, $a^{\log_bn}=n^{\log_ba}$, $n^{n/2}\leq n! \leq n^n$, $\log x^a = a\log x$
    \item $\log (a\times b) = \log a + \log b$, $\log(a/b) = \log a - \log b$
    \item $a + ar + ar^2 + ... + ar^{n-1} = \frac{a(r^n-1)}{r-1}$
    \item $1 + \frac{1}{2} + \frac{1}{2^2} + ... + \frac{1}{2^n} = 2(1-\frac{1}{2^{n+1}})$
    \item $1 + 2 + 4 + ... + 2^n = 2^{n+1}-1$
\end{itemize}


\hrule

\newpage

\begin{enumerate}
    \item \textbf{(10 pts)} \underline{Prefix-Sum Problem}: Given an array $A$ of integers, is there a sub-array of $A$ such that if we add all the elements of the sub-array, we get a sum of zero? By sub-array, we mean a contiguous sub-array. \\

    \textbf{Naive Brute Force Algorithm (Pseudocode)}:
    \begin{verbatim}
        for i = 1 to n {
            for j = i to n {
                sum = A[i] + A[i + 1] + ... + A[j]
                if sum == 0 return true
            }
        }
        return false
    \end{verbatim}
    Derive the runtime of the above algorithm in detail and determine its Big-O upper bound. You must show the derivation of the end result. Simply stating the final answer without any derivation steps will result in zero points.

    \ifans
    \color{blue}
    \color{black}
    \fi

    \hfill\break

    \item \textbf{(10 pts)} Given an array of integers, write an efficient algorithm to check whether there is a sub-array whose sum is a multiple of a given integer $K$. You must use a strategy of computing prefix-sums. Give a detailed analysis of the runtime or expected runtime of your algorithm. 
    
    \ifans
    \color{blue}
    \color{black}
    \fi

    \hfill\break

    \item \textbf{(10 pts)} Prove that the height of a heap tree is $\lfloor\log_2(n)\rfloor$ where $n$ is the number of elements in the tree.
    
    \ifans
    \color{blue}
    \color{black}
    \fi

    \hfill\break

    \item \textbf{(10 pts)}
    \begin{enumerate}
        \item \textbf{(5 pts)} Where in a max heap might the smallest element reside, assuming that all elements are distinct?
        
        \ifans
        \color{blue}
        \color{black}
        \fi
    
        \hfill\break

        \item \textbf{(5 pts)} Is an array that is sorted in increasing order a min heap? Justify your answer with an explanation.

        \ifans
        \color{blue}
        \color{black}
        \fi
    
        \hfill\break
    \end{enumerate}

    \hfill\break

    \item \textbf{(10 pts)} Consider an array $H$ of size $n \geq 2$, containing distinct positive integers and an integer $K$ such that $1 \leq K \leq n$. Without sorting the array, write an algorithm with input $H$ and $K$ to find the $K^{th}$ largest integer in $H$. Your algorithm must use a heap data structure. Write pseudocode for your algorithm and also present an analysis of its correctness and the runtime in terms of $K$ and $n$. 
\end{enumerate}

\end{document}




%%% Local Variables:
%%% mode: latex
%%% TeX-master: t
%%% End:
